%% ACM journal template — acmtog (single-column, no page limit).
%% Build: latexmk -pdf main.tex     (requires acmart.cls in TEXMF or same dir)

\documentclass[acmtog]{acmart}

%% --------- packages ---------
\usepackage{booktabs}
\usepackage{listings}
\usepackage{xcolor}
\usepackage{graphicx}
\usepackage{tikz}
\usepackage{enumitem}
\usepackage{tabularx}
\usepackage{float}
\usetikzlibrary{arrows.meta,positioning,shapes.geometric,fit,calc}

%% Allow line breaks in \texttt at underscores, dots, slashes, hyphens
\renewcommand\_{\textunderscore\allowbreak}
\newcommand{\tcode}[1]{{\small\texttt{#1}}}

\lstset{
  basicstyle=\footnotesize\ttfamily,
  breaklines=true,
  frame=single,
  xleftmargin=0pt,
  resetmargins=true,
  numbers=left,
  numberstyle=\tiny\color{gray},
}

%% Global: allow slightly looser line-breaking in two-column mode
\tolerance=1500
\emergencystretch=1em

%% --------- suppress ACM-specific fields (pre-submission draft) ---------
\setcopyright{none}
\settopmatter{printacmref=false, printfolios=true}
\renewcommand\footnotetextcopyrightpermission[1]{}
\acmJournal{JACM}
\acmVolume{1}
\acmNumber{1}
\acmArticle{1}
\acmYear{2026}
\acmMonth{4}
\copyrightyear{2026}
\acmDOI{}
\acmISBN{}
\acmPrice{}

\title{HLD-as-Source-of-Intent: A Ten-Agent Closed-Loop System for
  Intent-Based Networking Grounded in RFC~9315}

\author{Marc Bruy\`ere-Yamada}
\affiliation{%
  \institution{T\'el\'ecom SudParis}
  \city{\'Evry-Courcouronnes}
  \country{France}
}
\affiliation{%
  \institution{JFLI --- Japanese-French Laboratory for Informatics}
  \city{Tokyo}
  \country{Japan}
}
\email{marc@bruyere-yamada.name}

\author{Anonymous Co-authors}
\affiliation{%
  \institution{TBD}
  \country{}
}

\begin{abstract}
We present an Intent-Based Networking (IBN) closed-loop control system
in which the operator's High-Level Design (HLD) document, stored as
Markdown under version control, is the sole source of intent.
A~\texttt{git commit} on the HLD is the only external trigger: the
system reconciles live network state to the committed design through a
pipeline of ten specialised agents mapped to the stages defined in
RFC~9315, backed by a Neo4j graph Single Source of Truth (SSoT) with
nine ontology layers and five concurrent model states (What-If,
Candidate, Plan-of-Record, Deployed, As-Built). A three-tier memory
architecture separates structured state (Neo4j), working memory
(Live-Memory, MCP-attached), and long-term knowledge (Graph-Memory with
BGE-M3 vector embeddings and ontology-guided entity extraction).

We describe the two-loop control decomposition (autonomous inner loop,
human-in-the-loop outer loop), the nine-layer knowledge graph ontology
with its six cross-layer relationship types, the ten agents with their
precise inputs, outputs, dispatch tables, and write-layer assignments,
and the three-tier memory with its LLM-driven consolidation pipeline.
We report an end-to-end implementation covering multi-vendor rendering
across VyOS, Nokia SR~Linux, and FRRouting, shipped through six
incremental ``slices'' on a live Containerlab testbed with 16 nodes on a
12\,GB ARM64 host. Configuration is pushed through both
CLI-over-\texttt{docker exec} and schema-validated gNMI paths, with the
same A3-rendered intent artefact feeding either transport. The core
pipeline is deterministic and LLM-free; large language models are used
exclusively for memory consolidation and knowledge extraction, keeping
the fulfilment path's cycle time under 10\,seconds and its correctness
properties non-probabilistic.
\end{abstract}

\keywords{intent-based networking, RFC 9315, closed-loop control,
  GitOps, knowledge graphs, multi-agent systems, gNMI, YANG,
  network automation, Neo4j}

\begin{document}
\maketitle

%% ===================================================================
\section{Introduction}\label{sec:intro}
%% ===================================================================

\subsection{The HLD--Production Gap}

Modern enterprise and campus networks continue to be operated through
an uncomfortable split: architects produce dense High-Level Design
(HLD) documents that capture the \emph{intent} of the network---its
sites, zones, VLAN plans, firewall policies, routing design---while
operators, at a much later stage, translate those documents by hand
into device-specific configuration. The HLD is authoritative on paper
but inert in production: it does not execute, it does not verify
itself against the live topology, and it inevitably drifts from the
devices it is meant to describe.

This split has several well-documented consequences. First, the HLD
becomes a stale artefact within weeks of initial deployment: operators
make out-of-band changes to address urgent issues, and the document
is never backfilled. Second, the translation from HLD to device
configuration is manual, error-prone, and vendor-specific; an
architect's VLAN plan expressed in a prose table must be separately
re-encoded as VyOS \texttt{set} commands, Nokia SR~Linux \texttt{set /}
statements, or FRR \texttt{vtysh} directives. Third, the
\emph{rationale} for a design choice---captured naturally in prose and
margin notes in the HLD---is lost the moment the intent is translated
into a structured policy form. When a firewall rule is later found to
be overly permissive, the operator who must decide whether to narrow
it has no machine-accessible record of \emph{why} the rule was
written that way.

\subsection{Intent-Based Networking and RFC~9315}

Closing this gap is the central premise of Intent-Based Networking
(IBN), formalised by the IETF in RFC~9315~\cite{rfc9315}. The RFC
defines intent as an abstract, high-level policy that the network
infrastructure itself is responsible for rendering, deploying, and
continuously verifying through two nested feedback loops. The
\emph{fulfilment} phase ingests intent, decomposes it into
vendor-specific configuration, and orchestrates its deployment. The
\emph{assurance} phase monitors the deployed state, assesses
compliance against the intent-derived Plan-of-Record, and either
auto-remediates or escalates to a human operator. RFC~9315 further
distinguishes the \emph{inner loop} (autonomous, machine-to-machine,
seconds-to-minutes cadence) from the \emph{outer loop}
(human-in-the-loop, minutes-to-days cadence).

Despite a decade of vendor activity around IBN---Cisco DNA Center,
Juniper Apstra, Arista CloudVision, Nokia NSP---most deployed systems
remain tightly coupled to a single vendor's management plane and
treat intent as a GUI-authored artefact inside a controller. Operators
end up writing intent twice: once in the HLD (for humans) and once in
the controller (for machines). The two representations drift, and the
controller's intent model is a proprietary abstraction that cannot be
diffed, branched, or code-reviewed with standard development tooling.

\subsection{Contribution: HLD-as-Intent}

This paper describes a research prototype, developed on an open
Containerlab testbed~\cite{containerlab}, that treats the HLD Markdown
document itself as the source of intent. A~\texttt{git commit} on the
HLD is the only external trigger the system accepts: everything
downstream---parsing the change, updating the graph-structured Single
Source of Truth (SSoT), decomposing intent into policy, rendering
vendor-specific configuration, planning the migration with blast-radius
analysis, pushing to the devices, and verifying the outcome against the
committed design---is carried out autonomously by a pipeline of ten
specialised agents whose roles map directly onto the stages of the
RFC~9315 control loop. Architects edit the document they already know
how to write; the loop handles the rest.

The contributions of this paper are fivefold:

\begin{enumerate}[label=\textbf{C\arabic*.}]
  \item \textbf{HLD-as-Intent.} We demonstrate that a conventional
    Markdown HLD, augmented only with a structured Device Inventory
    Table and a VLAN assignment table, is sufficient to drive a full
    IBN fulfilment-and-assurance loop end-to-end, with \texttt{git}
    providing provenance, diffing, branching, and human review for
    free. The HLD is not a wrapper around a proprietary intent
    model; it \emph{is} the intent, and the system parses it
    directly using a 584-line Markdown parser with no LLM in the
    parsing path.

  \item \textbf{A ten-agent decomposition grounded in RFC~9315.} We
    map the abstract stages of RFC~9315~\S5 onto ten concrete agents
    (A1--A10) with narrow, testable responsibilities. Each agent
    writes to a well-defined subset of a nine-layer graph ontology
    (3,969~lines of Python across the ten agents), each emits
    structured working-memory notes that can be consolidated into
    long-term knowledge, and each publishes typed events to a Redis
    Streams event bus. The decomposition is fine-grained enough that
    adding a new vendor requires exactly one line per dispatch table
    plus a Jinja2 template chain---no agent code changes.

  \item \textbf{A nine-layer knowledge-graph ontology.} The Neo4j
    SSoT organises network state across nine layers
    (L1~Infrastructure through L9~People/Process) connected by six
    typed cross-layer relationships (\texttt{REALIZED\_BY},
    \texttt{ABSTRACTS}, \texttt{AGGREGATES}, \texttt{CONTAINS},
    \texttt{TRAVERSES}, \texttt{PRECEDED\_BY}). Five concurrent model
    states (\texttt{WHAT\_IF}, \texttt{CANDIDATE}, \texttt{POR},
    \texttt{DEPLOYED}, \texttt{AS\_BUILT}) co-exist as overlapping
    subgraphs, enforced by a controller with eight legal transition
    edges and four rollback paths. The ontology is codified in 616
    lines of Cypher DDL with profile-based constraint bundles.

  \item \textbf{A three-tier memory architecture.} Beyond the Neo4j
    SSoT (Tier~1: \emph{what is}) we layer a Live-Memory tier
    (Tier~2: \emph{why it is so}---working memory for agent
    reasoning, exposed via MCP Streamable HTTP) and a Graph-Memory
    tier (Tier~3: \emph{what was learned}---long-term knowledge graph
    with BGE-M3 1024-dimension embeddings in Qdrant, ontology-guided
    entity extraction across 14 entity classes and 8 relationship
    types). The three tiers are decoupled: the pipeline writes only
    to Tiers~1 and~2; consolidation from Tier~2 to Tier~3 runs
    asynchronously and is the only point where an LLM is invoked.

  \item \textbf{An incrementally shipped multi-vendor implementation.}
    We report on six ``slices'' of shipped functionality, culminating
    in a schema-validated gNMI push path for Nokia SR~Linux running
    side-by-side with CLI-over-\texttt{docker-exec} paths for VyOS
    and FRRouting~\cite{frr}. The testbed runs 16 Containerlab nodes
    (3~SR~Linux, 2~FRR, 4~VyOS, 7~host stubs) on a single 12\,GB
    ARM64 host with a full-cycle time of 8--10\,s. The implementation
    totals 140 unit tests across 23 test suites with zero regressions
    across slices.
\end{enumerate}

\subsection{Paper Organisation}

The remainder of this paper is organised as follows.
Section~\ref{sec:cl} describes the closed-loop architecture and its
grounding in RFC~9315, including the trigger-to-push pipeline, the
approval gate, and the invariants that make the system introspectable.
Section~\ref{sec:ssot} describes the Neo4j nine-layer ontology, the six
cross-layer relationship types, the five concurrent model states, and
the profile-based schema governance.
Section~\ref{sec:agents} details each of the ten agents, covering their
inputs, outputs, Neo4j write targets, dispatch tables, and key design
decisions.
Section~\ref{sec:memory} covers the three-tier memory architecture,
the consolidation pipeline, the IBN Lifecycle Ontology, and the
graph-guided RAG query path.
Section~\ref{sec:impl} reports on the incremental implementation across
six slices, the testbed configuration, measurements, and failure modes
encountered and resolved.
Section~\ref{sec:discuss} discusses related work and lessons learned.
Section~\ref{sec:conclusion} concludes.


%% ===================================================================
\section{The Closed Loop}\label{sec:cl}
%% ===================================================================

\subsection{Grounding in RFC~9315}\label{sec:cl:rfc}

RFC~9315~\cite{rfc9315} decomposes an IBN system into two nested
feedback loops. The \emph{inner loop} is autonomous: it runs on the
order of seconds to minutes, compares the network's observed
operational state against the intent-derived Plan-of-Record, detects
drift, classifies severity, and either self-corrects (re-render,
re-push) or escalates. The \emph{outer loop} is on human timescales:
operators review reports, amend the intent, and re-submit; the system
adapts accordingly. RFC~9315 further partitions the work into
\emph{fulfilment} (\S5.1: intent ingestion, policy translation, config
rendering, orchestration) and \emph{assurance} (\S5.2: monitoring,
assessment, remediation, abstraction, reporting).

Our system retains this two-loop, two-phase decomposition verbatim.
Table~\ref{tab:rfc-mapping} shows the mapping between RFC~9315
subsections and our ten agents. Beyond the standard decomposition, we
make three additional architectural commitments:

\begin{table}[t]
\small
\centering
\caption{Mapping of RFC~9315 stages to the ten agents.}
\label{tab:rfc-mapping}
\begin{tabular}{@{}llll@{}}
\toprule
\textbf{RFC~9315} & \textbf{Phase} & \textbf{Agent} & \textbf{Name}\\
\midrule
\S5.1.1 & Fulfilment & A1 & Ingestion\\
\S5.1.2 & Fulfilment & A2 & Intent$\rightarrow$Policy\\
\S5.1.2 & Fulfilment & A3 & Policy$\rightarrow$Config\\
\S5.1.2 & Fulfilment & A4 & Planning\\
\S5.1.3 & Fulfilment & A5 & Orchestration\\
\S5.2.1 & Assurance  & A6 & Monitoring\\
\S5.2.2 & Assurance  & A7 & Assessment\\
\S5.2.3 & Assurance  & A8 & Action\\
\S5.2   & Assurance  & A9 & Abstraction\\
\S5.2   & Assurance  & A10 & Reporting\\
\bottomrule
\end{tabular}
\end{table}

\begin{enumerate}[label=\textbf{AC\arabic*.}]
  \item \textbf{A single external trigger.} The only event that
    enters the system from outside is \texttt{git commit} on the HLD
    repository. A \texttt{post-commit} hook invokes the pipeline entry
    point \texttt{HldCommit\-Pipeline.run(\allowbreak commit, hld\_path)}, which is
    the sole supported way to start a fulfilment cycle. The hook
    runs the pipeline in a background process, emits a banner to the
    terminal (\texttt{[ibn-hook] HLD changed in <sha>}), and deposits
    a pending-approval marker when the approval gate fires. Everything
    else is internal. There is no GUI, no REST API, no controller; the
    intent surface is the text editor and the \texttt{git} CLI.

  \item \textbf{A single human checkpoint.} The inner loop is
    autonomous by design, but change-of-intent operations (new
    devices, decommissions, policy scope changes) are gated by an
    explicit approval step between Agent~4 (Planning) and Agent~5
    (Orchestration). The approval artefact is a Markdown file
    deposited under \texttt{.git/\allowbreak{}ibn-pending-approval/\allowbreak{}<sha>.md}
    containing the severity classification, blast radius, affected
    devices, and a forward/inverse summary of the planned change. An
    operator approves by running \texttt{python -m ibn.tools.approve
    <sha>} or rejects with \texttt{python -m ibn.tools.approve reject
    <sha> --reason "..."}. This preserves RFC~9315's outer-loop human
    oversight but removes the GUI-controller dependency.

  \item \textbf{Deterministic pipeline, LLM-assisted memory.} The
    core pipeline (A1$\rightarrow$A5 fulfilment, A6$\rightarrow$A8
    assurance) is entirely LLM-free. LLMs enter the architecture
    exclusively at the memory consolidation and knowledge extraction
    boundary (Section~\ref{sec:memory}), which runs asynchronously
    after cycle completion. This separation is intentional: the
    pipeline's correctness properties---idempotent renders, atomic
    gNMI commits, rollback chains via \texttt{PRECEDED\_BY}---depend
    on determinism. The memory tier's value proposition---turning
    thousands of structured agent notes into queryable operational
    knowledge---is precisely where LLM summarisation excels. The
    architecture puts each capability where it pays off.
\end{enumerate}


\subsection{Trigger-to-Push Pipeline}\label{sec:cl:pipeline}

When a \texttt{git commit} touches the HLD file
(\texttt{Enterprise\_Campus\_\allowbreak{}Network\_HLD.md}), the \texttt{post-commit}
hook spawns \texttt{HldCommit\-Pipeline.run(\allowbreak commit\_sha,\allowbreak{} hld\_path)}.
The pipeline carries out twelve stages in order; each stage is
described below with its agent, inputs, outputs, and the Neo4j layer
it writes to.

\begin{enumerate}[label=\textbf{Stage~\arabic*.},leftmargin=3em]
  \item \textbf{Diff (parser).}
    The HLD Markdown is parsed by a deterministic 584-line parser
    (\texttt{hld\_parser.py}) that extracts structured sections:
    the Device Inventory Table, population VLAN tables, DMZ VLAN
    tables, and firewall policy matrices. The parser produces a
    typed \texttt{HldChangeset} dataclass with fields:
    \texttt{populations\_added}, \texttt{populations\_removed},
    \texttt{populations\_modified}, \texttt{dmz\_added},
    \texttt{dmz\_removed}, \texttt{dmz\_modified},
    \texttt{devices\_added}, \texttt{devices\_removed},
    \texttt{devices\_modified}. Each field is a list of
    domain-specific typed records (e.g.\ \texttt{HldDevice} with
    fields \texttt{hostname}, \texttt{vendor}, \texttt{platform},
    \texttt{role}, \texttt{mgmt\_ipv4}). The changeset is
    computed by diffing the parsed state of the new commit against
    the previous commit's parsed state.

  \item \textbf{Agent~1 (Ingestion).}
    The ingestion agent (\texttt{agent1\_ingestion.py}, 584~lines)
    receives the changeset and upserts nodes into Neo4j at layers
    L1 (Device, VLAN, Site) and L4 (Intent). Every node is created
    with \texttt{modelState = CANDIDATE}. Idempotency is enforced
    by keying each intent on its HLD origin:
    \texttt{HLD:<path>:<line\_number>}. If an intent with the same
    origin already exists, the agent updates rather than duplicates.
    Conflict detection queries all active intents and checks for
    \texttt{CONFLICTING\_ACTION} (same target, incompatible action)
    or \texttt{DUPLICATE} (same target, same action). Each write
    is dual-sinked: a \texttt{live\_note()} call to the
    \texttt{ibn-candidate-<version>} Live-Memory space records the
    write's context (what was written, what was considered, why).
    Events are published to the Redis Streams event bus:
    \texttt{intent.ingested}, \texttt{device.added},
    \texttt{device.modified}, \texttt{device.removed}.

    For device changes, Agent~1 also handles lifecycle state: new
    devices enter as \texttt{lifecycleState = PLANNED}, modified
    devices are updated in place, and removed devices are marked
    \texttt{RETIRED} (soft delete---the node stays in the graph for
    audit).

  \item \textbf{Provisioning.}
    For any device deltas, the Provisioner
    (\texttt{provisioner.py}) mutates the Containerlab topology YAML
    and triggers a \texttt{clab deploy --reconfigure}. The
    Provisioner has two code paths:

    \begin{itemize}
      \item \texttt{provision(device\_entry)}: single-device path.
        Adds the node to the YAML, runs \texttt{clab deploy
        --reconfigure}, waits for the container, records a
        \texttt{LifecycleEvent} at L8.
      \item \texttt{provision\_batch(device\_entries)}: multi-device
        path (used when $|\text{devices\_added}| > 1$). Performs a
        single YAML write and a single \texttt{clab deploy
        --reconfigure} for the entire batch. This was a mid-Slice-5
        fix: the per-device path ran \texttt{clab deploy
        --reconfigure} once per device, and since \texttt{--reconfigure}
        re-processes the entire topology each time, four sequential
        deploys exceeded the 300\,s timeout on the third iteration.
        The batch path completes in $\sim$10\,s.
    \end{itemize}

    The vendor$\rightarrow$containerlab-kind mapping is driven by a
    dispatch table \texttt{\_KIND\_MAP}:

    \begin{lstlisting}[language=Python,caption={Vendor dispatch for Containerlab kinds.}]
_KIND_MAP = {
  ("nokia","srlinux"): {"kind":"nokia_srlinux",
    "image":"ghcr.io/nokia/srlinux:latest",
    "type":"ixrd2l"},
  ("vyos","vyos"): {"kind":"linux",
    "image":"ghcr.io/sysoleg/vyos-container:latest"},
  ("frr","frr"): {"kind":"linux",
    "image":"ibn-frr:local"},
}
    \end{lstlisting}

    Adding a new vendor to the provisioning path is one entry in
    this table plus a Dockerfile for the container image.

  \item \textbf{Agent~2 (Intent$\rightarrow$Policy).}
    The policy decomposition agent (\texttt{agent2\_intent\_policy.py},
    221~lines) receives the list of intent IDs produced by Stage~2 and
    decomposes each into \texttt{Policy} and \texttt{FirewallRule}
    nodes at L4. Policy IDs follow the pattern
    \texttt{POL-<intent\_id>}; rule IDs follow
    \texttt{RUL-<intent\_id>-<idx:02d>}. The decomposition is driven
    by an \texttt{ACCESS\_MATRIX} dispatch table that maps a target
    string (e.g.\ \texttt{"Internet+DMZ"}, \texttt{"DNS-DMZ"}) to a
    list of concrete firewall rules with fields \texttt{destZone},
    \texttt{action}, \texttt{protocol}, \texttt{destPort}, and
    \texttt{description}.

    Zone inference is currently rule-based: DMZ targets map to
    \texttt{DMZ\_INFRA}, guest/contractor/visitor targets to
    \texttt{GUEST}, all others to \texttt{USER}. Device-only commits
    (no population or DMZ changes) short-circuit this stage entirely:
    Agent~2 reports \texttt{skipped} and the pipeline continues with
    the existing intent set.

  \item \textbf{Agent~4 (Planning).}
    The planning agent (\texttt{agent4\_planning.py}, 243~lines)
    produces a \texttt{MigrationPlan} node at L7 with:
    \begin{itemize}
      \item A \textbf{severity classification} computed by a
        deterministic heuristic:
        \texttt{HIGH} if devices were added or removed;
        \texttt{MEDIUM} if populations or DMZ segments were removed;
        \texttt{LOW} if populations or DMZ segments were added or
        modified; \texttt{INFO} if no changes detected.
      \item A \textbf{blast radius}: the count and list of affected
        devices, computed by traversing
        \texttt{Configuration}$\rightarrow$\texttt{Device} edges
        (\texttt{APPLIES\_TO}).
      \item A \textbf{forward summary} (what will change) and an
        \textbf{inverse summary} (how to undo it), both in
        human-readable prose.
    \end{itemize}
    Plan IDs follow the pattern \texttt{PLAN-<commit\_sha[:8]>}.
    If a prior \texttt{PENDING} plan exists for the same intent set,
    it is automatically marked \texttt{SUPERSEDED} before the new
    plan is written.

  \item \textbf{Approval gate.}
    If severity is \texttt{HIGH} (device additions or removals), the
    pipeline deposits a Markdown summary under
    \texttt{.git/\allowbreak{}ibn-pending-approval/\allowbreak{}<sha>.md} and halts. The
    summary contains: severity, blast radius, affected device list,
    forward and inverse summaries, and the command to approve. The
    operator reviews and runs:

    \begin{lstlisting}[language=bash,numbers=none]
python -m ibn.tools.approve <sha>
    \end{lstlisting}

    The approve tool (\texttt{ibn.tools.approve}) transitions the
    \texttt{MigrationPlan} from \texttt{PENDING} to \texttt{APPROVED},
    records the operator identity and timestamp, cleans up the pending
    Markdown, and invokes
    \texttt{HldCommit\-Pipeline.\allowbreak{}resume\_from\_approval(\allowbreak{}plan\_id)}.

    For \texttt{LOW}/\texttt{INFO} severity, the gate auto-approves
    (configurable via \texttt{IBN\_AUTO\_APPROVE=1} for test
    environments).

  \item \textbf{Agent~3 (Policy$\rightarrow$Config).}
    The config rendering agent (\texttt{agent3\_policy\_config.py},
    420~lines) is the most data-intensive stage. It executes ten
    named Cypher queries against Neo4j to build a rendering context:
    \texttt{firewalls} (FirewallPair$\rightarrow$Device with
    \texttt{modelState IN ['POR','DEPLOYED']}), \texttt{zones}
    (ordered by securityLevel), \texttt{segments}
    (Segment$\rightarrow$VLAN via \texttt{MAPPED\_TO\_VLAN}),
    \texttt{policies} (FirewallRule with legacy field coalescing),
    \texttt{intents}, \texttt{fw\_interfaces},
    \texttt{nat\_rules}, \texttt{fw\_transit} (VLAN 300 with
    subnets), \texttt{switches} (non-firewall devices excluding
    \texttt{RETIRED}), and \texttt{vlans}.

    For each device, Agent~3 resolves the appropriate Jinja2
    template chain from the \texttt{\_TEMPLATE\_CHAINS} dispatch
    table:

    \begin{lstlisting}[language=Python,caption={Template chains per vendor.}]
_TEMPLATE_CHAINS = {
  "vyos":    ["vyos_base.j2", "vyos_policies.j2",
              "vyos_nat.j2", "vyos_ha.j2"],
  "srlinux": ["srl_base.j2", "srl_interfaces.j2",
              "srl_vlans.j2"],
  "frr":     ["frr_base.j2", "frr_interfaces.j2",
              "frr_routing.j2"],
}
    \end{lstlisting}

    Each template in the chain is rendered against the Neo4j context
    and concatenated. The resulting configuration blob is written as a
    \texttt{Configuration} node at L5 with a content hash
    (\texttt{SHA-256[:16]}) for change detection. A diff against the
    previous configuration for the same device is computed to
    determine whether the render produced net-new content.

  \item \textbf{Agent~5 (Orchestration).}
    The orchestration agent (\texttt{agent5\_orchestration.py},
    $\sim$350~lines) pushes each rendered configuration to its live
    device. The executor is resolved from a two-level dispatch:

    \begin{lstlisting}[language=Python,caption={Executor dispatch with transport override.}]
_VENDOR_EXECUTORS = {
  "vyos":    _default_ssh_executor,
  "srlinux": _srlinux_ssh_executor,
  "frr":     _frr_executor,
}

def _resolve_executor(platform):
  if platform == "srlinux" and
     os.environ.get("IBN_SRLINUX_TRANSPORT")
     == "gnmi":
    return _srlinux_gnmi_executor
  return _VENDOR_EXECUTORS.get(platform,
    _default_ssh_executor)
    \end{lstlisting}

    Four executor implementations exist:
    \begin{description}[style=sameline,leftmargin=2em]
      \item[VyOS:] \texttt{paramiko} SSH to the device, pipe
        configuration commands.
      \item[SR~Linux (CLI):] \texttt{docker exec -i <container>
        sr\_cli}, wrapping commands in
        \texttt{enter candidate / ... / commit now / quit}.
      \item[SR~Linux (gNMI):] A narrow CLI$\rightarrow$YANG path
        translator (\texttt{\_srlinux\_yang.py}, $\sim$120~lines)
        converts the rendered \texttt{set /} commands into gNMI
        \texttt{(path, value)} tuples. The executor issues a single
        \texttt{SetRequest} via \texttt{pygnmi}~\cite{pygnmi} to
        \texttt{<mgmt-ipv4>:57400} with TLS (self-signed cert,
        \texttt{skip\_verify=True}). The translator has ten rules
        covering the exact line shapes the three \texttt{srl\_*.j2}
        templates emit; unknown lines raise
        \texttt{UnknownCliLine} to prevent silent config loss.
        Multiple leaves under the same container path are merged
        into a single gNMI update, reducing wire operations from
        10 to 7 for a typical render.
      \item[FRR:] \texttt{docker exec -i <container> vtysh},
        wrapping in \texttt{enable / configure terminal / ... /
        end / write memory}.
    \end{description}

    Each successful push records a \texttt{DeploymentEvent} node at
    L5. Devices with empty rendered content (e.g.\ device-only commits
    where templates produce no new output) are skipped.

  \item \textbf{Agent~6 (Monitoring).}
    The monitoring agent (\texttt{agent6\_monitoring.py}) collects
    post-deployment telemetry via CLI polling (\texttt{show}
    commands over SSH or \texttt{docker exec}). For each device it
    creates \texttt{Telemetry} nodes and upserts
    \texttt{OperationalState} nodes at L5 with
    \texttt{modelState = AS\_BUILT}. Metrics include interface
    state, VRRP role, BGP peer status, and OSPF neighbour count.
    The monitoring interval is configurable (default 30\,s for the
    inner loop).

  \item \textbf{Agent~7 (Assessment).}
    The assessment agent (\texttt{agent7\_assessment.py}, 417~lines)
    compares the \texttt{POR} and \texttt{AS\_BUILT} subgraphs. It
    fetches per-intent POR configurations and per-device As-Built
    state, then runs three registered drift detectors:

    \begin{description}[style=sameline,leftmargin=2em]
      \item[\texttt{firewall\_rule}:] Detects missing or extra rules
        by comparing POR FirewallRule sets against As-Built running
        configuration snippets.
      \item[\texttt{interface\_state}:] Detects interface
        admin/oper-state mismatches.
      \item[\texttt{vrrp\_role}:] Detects VRRP role mismatches
        (primary/backup).
    \end{description}

    Each drift produces a \texttt{ComplianceAssessment} node at L6
    with a verdict (\texttt{COMPLIANT}, \texttt{DEGRADED},
    \texttt{NON\_COMPLIANT}, \texttt{UNKNOWN}) and a compliance
    percentage. For non-compliant verdicts, a root-cause tracer
    traverses the dependency graph (via
    \texttt{REALIZED\_BY}/\allowbreak\texttt{AGGREGATES} edges) to identify
    the causal chain and assigns a confidence level
    (\texttt{HIGH} or \texttt{LOW}).

    If drift events are detected, an \texttt{Incident} node is
    created at L6, linked to the assessment. Severity classification:
    \texttt{CRITICAL} ($\geq$3 drift events), \texttt{HIGH}
    (non-compliant verdict), \texttt{MEDIUM} (default).

  \item \textbf{Agent~8 (Action).}
    The action agent (\texttt{agent8\_action.py}, 405~lines) receives
    incident data from Agent~7 and applies a five-level autonomy
    decision matrix:

    \begin{table}[H]
    \small
    \centering
    \caption{Agent~8 autonomy decision matrix.}
    \label{tab:autonomy}
    \begin{tabular}{@{}llcc@{}}
    \toprule
    \textbf{Severity} & \textbf{Action} & \textbf{Auto} & \textbf{Human}\\
    \midrule
    \texttt{INFO}     & Log only             & \checkmark & \\
    \texttt{LOW}      & Auto-remediate       & \checkmark & \\
    \texttt{MEDIUM}   & Remediate + verify   & \checkmark & \\
    \texttt{HIGH}     & Escalate             &            & \checkmark\\
    \texttt{CRITICAL} & Rollback + escalate  &            & \checkmark\\
    \bottomrule
    \end{tabular}
    \end{table}

    Auto-remediation re-invokes A3$\rightarrow$A5 on the narrowed
    artefact set for the affected devices. Rollback walks the
    \texttt{PRECEDED\_BY} chain one step to retrieve the previous
    \texttt{Configuration} node and re-pushes it. Every action is
    recorded as a \texttt{Remediation} node at L6 with fields:
    \texttt{remediationId}, \texttt{intentId}, \texttt{action},
    \texttt{success}, \texttt{error}, \texttt{escalated},
    \texttt{rollbackCount}.

  \item \textbf{Agents~9--10 (Abstraction, Reporting).}
    Agent~9 aggregates device-level observations into intent-level
    compliance scores using \texttt{ABSTRACTS} edge traversals.
    Agent~10 generates human-facing compliance reports. Both are
    designed but not yet fully implemented in the current prototype;
    their stubs emit aggregate metrics to the event bus and write
    summary data to the graph.
\end{enumerate}

\subsection{Pipeline Timing}

A typical no-op HLD edit (comment change, no structural delta)
completes the full twelve-stage cycle in 8--10\,s against the current
testbed (16 Containerlab nodes on a 12\,GB ARM64 host). The
dominant costs are Neo4j Bolt round-trips ($\sim$200\,ms per Cypher
query, ten queries in A3) and container-exec I/O ($\sim$1\,s per
device push). A device-addition commit, which goes through the
approval gate and a \texttt{clab deploy --reconfigure}, completes
in 20--40\,s of wall-clock time once the human approval is granted.
The approval itself is out-of-band and untimed.


\subsection{Invariants}\label{sec:cl:invariants}

Three invariants hold across every stage of the pipeline and are
load-bearing for the system's correctness:

\begin{description}
  \item[I1: Every agent write is dual-sink.] Every node or edge an
    agent writes to Neo4j is accompanied by a structured
    \texttt{live\_note()} call into a Live-Memory space keyed on the
    current model state. This preserves the \emph{why} of the write
    (prompt, inputs considered, alternatives rejected) alongside the
    \emph{what} (the graph mutation itself). Invariant~I1 is enforced
    by the \texttt{BaseAgent} class: every agent inherits from
    \texttt{BaseAgent}, which provides a \texttt{\_note()} helper
    that writes to both sinks atomically. An agent that calls
    \texttt{neo4j.create\_*} without a corresponding
    \texttt{\_note()} violates the contract; this is detectable by
    static analysis (grep for unmatched writes).

  \item[I2: Every node carries \texttt{modelState}.] The five-valued
    \texttt{modelState} property (\texttt{WHAT\_IF}, \texttt{CANDIDATE},
    \texttt{POR}, \texttt{DEPLOYED}, \texttt{AS\_BUILT}) is written on
    every node creation and transitioned through well-defined edges. A
    legal transition table (eight edges, four rollback paths) is
    enforced by the \texttt{ModelStateController}; illegal transitions
    raise \texttt{ModelStateTransitionError} at write time rather than
    corrupting the graph silently. Every transition creates a
    \texttt{LifecycleTransition} node at L8 linked to the affected
    intent via a \texttt{HAS\_TRANSITION} edge.

  \item[I3: Every agent execution is recorded.] Each agent run
    produces an \texttt{AgentExecution} node at L9 with the agent
    ID, start/end timestamps, inputs summary, outputs summary,
    success/failure status, and the intent ID being processed. This
    provides full traceability from intent to device to compliance
    assessment and back.
\end{description}

Together, I1--I3 ensure the system is fully introspectable: at any
point, an operator can ask ``what is POR right now?'' (query
\texttt{modelState = 'POR'} in Neo4j), ``what changed in the last
cycle?'' (query recent \texttt{AgentExecution} and
\texttt{LifecycleTransition} nodes), and ``why did the agent do
this?'' (read the corresponding Live-Memory note) and get machine-
structured answers.


%% ===================================================================
\section{SSoT: Nine-Layer Knowledge-Graph Ontology}\label{sec:ssot}
%% ===================================================================

The Single Source of Truth (SSoT) is a Neo4j~\cite{neo4j} property
graph whose schema is organised as a \emph{layered ontology}. The
ontology is codified in 616~lines of Cypher DDL
(\texttt{Neo4j\_Schema\_Cypher.cypher}) with constraints, indexes,
and seed-data templates. Layering is load-bearing: each agent in the
pipeline of Section~\ref{sec:agents} writes into a small,
well-defined subset of layers, and every cross-layer relationship
has an explicit, audit-friendly semantics. The ontology is inspired
by Google's MALT~\cite{malt} multi-level topology abstraction, but
specialised for the IBN lifecycle rather than for datacentre traffic
engineering, and aligned with the fulfilment/assurance split of
RFC~9315~\cite{rfc9315}.


\subsection{The Nine Layers}\label{sec:ssot:layers}

Table~\ref{tab:layers} lists the nine layers with their canonical
node labels. Labels are reserved tokens: every node in the graph
carries exactly one \emph{primary} layer label (e.g.\
\texttt{Device}, \texttt{Intent}, \texttt{MigrationPlan}) and may
carry additional tag labels (e.g.\ \texttt{:Active}, \texttt{:POR}).

\begin{table*}[t]
\small
\centering
\caption{Nine-layer ontology: one canonical label set per layer, with the
agent(s) that write to each layer.}
\label{tab:layers}
\begin{tabularx}{\textwidth}{@{}l l X l@{}}
\toprule
\textbf{Layer} & \textbf{Name} & \textbf{Representative labels} & \textbf{Writers}\\
\midrule
L1 & Infrastructure &
  \texttt{Site}, \texttt{Device}, \texttt{Interface},
  \texttt{PhysicalLink}, \texttt{LogicalLink}, \texttt{VLAN},
  \texttt{VRF}, \texttt{Subnet}, \texttt{IPAddress} &
  A1\\
L2 & Topology &
  \texttt{Topology}, \texttt{TopologyLayer}, \texttt{Zone},
  \texttt{Segment}, \texttt{FirewallPair}, \texttt{StackGroup} &
  A1 (seed)\\
L3 & Architecture &
  \texttt{Architecture}, \texttt{ArchitectureVersion},
  \texttt{HLD}, \texttt{ARD}, \texttt{BusinessUseCase} &
  A1 (seed)\\
L4 & Policy \& Intent &
  \texttt{Intent}, \texttt{Policy}, \texttt{FirewallRule},
  \texttt{SGT}, \texttt{QoSPolicy} &
  A1, A2\\
L5 & Config \& State &
  \texttt{Configuration}, \texttt{DeploymentEvent},
  \texttt{Telemetry}, \texttt{OperationalState} &
  A3, A5, A6\\
L6 & Incidents &
  \texttt{Alert}, \texttt{Incident}, \texttt{RootCause},
  \texttt{ComplianceAssessment}, \texttt{Remediation} &
  A7, A8\\
L7 & Change Mgmt &
  \texttt{ChangeRequest}, \texttt{MigrationPlan} &
  A4\\
L8 & Lifecycle &
  \texttt{LifecyclePhase}, \texttt{LifecycleTransition},
  \texttt{LifecycleEvent} &
  Provisioner, MSC\\
L9 & People/Process &
  \texttt{Operator}, \texttt{Team}, \texttt{Ticket},
  \texttt{SLA}, \texttt{AgentExecution} &
  All agents\\
\bottomrule
\end{tabularx}
\end{table*}

The layering is strictly semantic, not a storage partition: Neo4j
holds all labels in a single graph so that queries can traverse
across layers in a single Cypher statement. The value of the layer
is that it constrains \emph{who writes what}: Agent~1 only writes
L1/L4; Agent~3 only writes L5; Agent~7 only writes L6. This makes
each agent's blast radius on the graph bounded and testable. A
table of agents and their write layers is included in
Table~\ref{tab:layers}.

\paragraph{Layer design rationale.}
The nine layers correspond to nine qualitatively different concerns
in the network lifecycle. Layers L1--L3 capture the static
structure (infrastructure, topology, architecture) that changes on
the timescale of HLD revisions. L4 captures intent and policy,
which change on the timescale of business requirements. L5 captures
the dynamic configuration and operational state that changes on the
timescale of pipeline cycles. L6 captures incidents and compliance,
which are generated as side effects of assurance. L7 captures
planned changes, which exist between intent and deployment. L8
captures lifecycle transitions, which are the audit trail of state
evolution. L9 captures the human and agent actors involved.


\subsection{Cross-Layer Relationships}\label{sec:ssot:rels}

Six cross-layer relationship types form the backbone of the
ontology. They are the \emph{only} edges permitted to cross a layer
boundary; intra-layer edges (e.g.\ \texttt{CONNECTS\_TO} between two
\texttt{Interface} nodes, \texttt{PEER\_OF} between two
\texttt{Device} nodes) are free-form within a layer.

\begin{description}
  \item[\texttt{REALIZED\_BY} (descending realisation).]
    An \texttt{Intent} (L4) is \texttt{REALIZED\_BY} a set of
    \texttt{Policy}/\texttt{FirewallRule} nodes (L4), which are in
    turn \texttt{REALIZED\_BY} \texttt{Configuration} nodes (L5),
    which are \texttt{REALIZED\_BY} the live
    \texttt{OperationalState} (L5). The \texttt{REALIZED\_BY} chain
    is the primary audit trail: following it downward from an intent
    answers ``does this intent actually exist on the wire?''. Agent~7
    performs exactly this traversal during drift detection: it
    fetches the \texttt{POR} configuration for an intent via
    \texttt{REALIZED\_BY} edges and compares it against the
    \texttt{AS\_BUILT} operational state.

  \item[\texttt{ABSTRACTS} (ascending abstraction).]
    The inverse of \texttt{REALIZED\_BY}. Agent~9 uses
    \texttt{ABSTRACTS} edges to roll up device-level observations
    into intent-level compliance scores and then into site-level
    summaries.

  \item[\texttt{AGGREGATES} (set membership).]
    Distinct from realisation. A \texttt{FirewallPair} (L2)
    \texttt{AGGREGATES} two \texttt{Device} nodes (L1); a
    \texttt{Segment} (L2) \texttt{AGGREGATES} a set of
    \texttt{VLAN} nodes (L1). \texttt{AGGREGATES} is structural---it
    says ``these nodes form a group''---while \texttt{REALIZED\_BY}
    is functional---it says ``this node is the concrete form of that
    abstract node''.

  \item[\texttt{CONTAINS} (structural nesting).]
    A \texttt{Site} (L1) \texttt{CONTAINS} \texttt{Device} nodes;
    an \texttt{Architecture} (L3) \texttt{CONTAINS}
    \texttt{ArchitectureVersion} nodes; an \texttt{HLD} (L3)
    \texttt{CONTAINS} \texttt{Intent} nodes extracted from its
    sections. \texttt{CONTAINS} always points from a parent to its
    children within a single ownership hierarchy.

  \item[\texttt{TRAVERSES} (policy crossing).]
    A \texttt{FirewallRule} (L4) \texttt{TRAVERSES} a
    \texttt{FirewallPair} (L2). An \texttt{Intent} (L4)
    \texttt{TRAVERSES} a \texttt{Zone} (L2) when its scope is
    zone-level. Blast-radius analysis in Agent~4 is largely a
    \texttt{TRAVERSES} traversal: ``how many zones does this intent
    cross? how many firewall pairs does this rule set touch?''.

  \item[\texttt{PRECEDED\_BY} (temporal ordering).]
    Every \texttt{Configuration} node carries a
    \texttt{PRECEDED\_BY} edge to the configuration it replaced.
    Rollback in Agent~8 is implemented by walking
    \texttt{PRECEDED\_BY} one step and re-pushing the predecessor.
    The chain is unbounded; an operator can walk it back to the
    initial configuration to see the full history of a device's
    config evolution.
\end{description}


\subsection{Five Concurrent Model States}\label{sec:ssot:states}

Every node in the graph carries a \texttt{modelState} property
drawn from the set
$\mathcal{S} = \{$\texttt{WHAT\_IF}, \texttt{CANDIDATE},
\texttt{POR}, \texttt{DEPLOYED}, \texttt{AS\_BUILT}$\}$.
The five states co-exist \emph{simultaneously} as overlapping
subgraphs over the same node labels and relationship types;
selecting a state is a Cypher filter on \texttt{modelState}.
This design is directly inspired by MALT's~\cite{malt} observation
that multiple levels of abstraction must be maintained
simultaneously for a topology to be operationally useful.

\begin{description}
  \item[\texttt{WHAT\_IF}.]
    Exploratory scratch space used by Agent~4 during planning.
    Blast-radius simulations clone the affected subgraph into
    \texttt{WHAT\_IF} state, mutate it, and discard it after
    analysis. The production graph is never touched.

  \item[\texttt{CANDIDATE}.]
    State produced by a freshly committed HLD that has not yet been
    approved. Agent~1 writes new nodes directly as
    \texttt{CANDIDATE}. The \texttt{CANDIDATE} subgraph is the
    ``proposed future'': it co-exists with the current \texttt{POR}
    until the approval gate promotes it.

  \item[\texttt{POR} (Plan-of-Record).]
    The canonical intent: what the network \emph{should} be.
    Transition from \texttt{CANDIDATE} to \texttt{POR} happens only
    after the approval gate is cleared. All rendering queries in
    Agent~3 filter for \texttt{modelState IN ['POR', 'DEPLOYED']} to
    ensure they operate against approved state.

  \item[\texttt{DEPLOYED}.]
    Reflects what Agent~5 actually pushed, device by device. A
    \texttt{Configuration} node becomes \texttt{DEPLOYED} when the
    executor returns success. \texttt{DEPLOYED} may temporarily
    differ from \texttt{POR} during a rolling deployment when some
    devices have been updated and others have not.

  \item[\texttt{AS\_BUILT}.]
    The ground-truth view built from Agent~6's telemetry. May
    diverge from \texttt{DEPLOYED} in the presence of out-of-band
    changes, partial push failures, or device-side configuration
    mutations. Agent~7's drift detection is precisely a
    \texttt{POR} vs.\ \texttt{AS\_BUILT} set-difference over the
    subgraph.
\end{description}

\paragraph{Legal transitions.}
The \texttt{ModelStateController} (298~lines) enforces eight legal
transitions stored as a frozen set of \texttt{(from, to)} pairs:

\medskip
\noindent\textit{Forward:}
\texttt{WHAT\_IF}~$\rightarrow$~\texttt{CANDIDATE}
$\rightarrow$~\texttt{POR}
$\rightarrow$~\texttt{DEPLOYED}
$\rightarrow$~\texttt{AS\_BUILT}

\noindent\textit{Rollback:}
\texttt{CANDIDATE}~$\rightarrow$~\texttt{WHAT\_IF},
\texttt{POR}~$\rightarrow$~\texttt{CANDIDATE},
\texttt{DEPLOYED}~$\rightarrow$~\texttt{POR},
\texttt{AS\_BUILT}~$\rightarrow$~\texttt{DEPLOYED}
\medskip

Any other transition raises \texttt{ModelStateTransitionError}.
Every transition creates a \texttt{LifecycleTransition} node at L8
with fields: \texttt{transitionId}, \texttt{intentId},
\texttt{fromState}, \texttt{toState}, \texttt{reason},
\texttt{operator}, \texttt{timestamp}. The controller also assigns
each transition to a Live-Memory space via a prefix map:

\begin{lstlisting}[language=Python,caption={Live-Memory space prefix per model state.}]
_SPACE_PREFIXES = {
  "WHAT_IF":   "ibn-whatif",
  "CANDIDATE": "ibn-candidate",
  "POR":       "ibn-por",
  "DEPLOYED":  "ibn-deploy",
  "AS_BUILT":  "ibn-asbuilt",
}
\end{lstlisting}

Convenience methods \texttt{approve()}, \texttt{deploy()},
\texttt{verify()}, and \texttt{rollback()} wrap the generic
\texttt{transition()} call with the appropriate \texttt{(from, to)}
pair and a default reason string.


\subsection{Profiles and Schema Governance}\label{sec:ssot:profiles}

Neo4j constraints encode the ontology machine-checkably. Beyond
unique-key constraints on node identifiers (e.g.\
\texttt{CREATE CONSTRAINT FOR (d:Device) REQUIRE d.deviceId IS
UNIQUE}), the schema declares \emph{profiles}: named bundles of
Cypher constraints that capture deployment-specific invariants.
Three profiles are defined:

\begin{description}
  \item[Campus profile.] Every \texttt{Device} must be
    \texttt{CONTAINS}-linked to exactly one \texttt{Site}. Every
    \texttt{FirewallRule} must \texttt{TRAVERSES} exactly one
    \texttt{FirewallPair}. Every \texttt{VLAN} must have a
    non-null \texttt{vlanId}. This profile matches the reference
    HLD topology.

  \item[Branch profile.] Relaxes the Campus profile to allow
    devices without a \texttt{FirewallPair} (branch sites may
    use a cloud firewall), and allows VLANs without a local
    \texttt{Segment} assignment.

  \item[Lab profile.] The most permissive profile. Allows
    devices without site assignment (for ad-hoc containerlab
    testing) and relaxes firewall-pair constraints.
\end{description}

Violating a profile constraint is a deployment-time Cypher error,
not a runtime surprise. The active profile is set per environment
via the \texttt{IBN\_SCHEMA\_PROFILE} variable.

Schema evolution is versioned semantically
(\texttt{MAJOR.MINOR.PATCH}) with an explicit
\texttt{OntologyVersion} node tracking which version is active. A
six-month stability window on the \texttt{MAJOR} version gives
agents a contract they can rely on, and a named stored-procedure
library over canonical Cypher queries insulates agents from
\texttt{MINOR}-version column renames.


\subsection{Why Nine Layers, Not Fewer}\label{sec:ssot:why}

Early iterations of the ontology collapsed L7 (Change Management)
into L8 (Lifecycle) and L6 (Incidents) into L5 (Config \& State).
Both collapses damaged query ergonomics and introduced write-lock
contention:

\begin{itemize}
  \item Separating L6 from L5 lets the assurance agents (A7, A8)
    own their own layer without competing for write-locks with the
    fulfilment agents (A3, A5, A6) operating on L5. In a concurrent
    inner loop where A6 is writing telemetry to L5 while A7 is
    writing assessments, the two agents touch disjoint label sets.

  \item Separating L7 from L8 decouples \emph{planned} change
    (\texttt{MigrationPlan}, \texttt{ChangeRequest}) from
    \emph{actual} transitions (\texttt{LifecycleEvent},
    \texttt{LifecyclePhase}). This matters because the former are
    reasoned about \emph{prospectively} (Agent~4 creates them before
    approval) and the latter are reasoned about
    \emph{retrospectively} (the audit trail after the fact).

  \item Separating L9 from L7/L8 gives the agent execution audit
    trail its own namespace. Every agent writes an
    \texttt{AgentExecution} node to L9 on every run; mixing these
    high-frequency writes with the lower-frequency
    \texttt{MigrationPlan} nodes of L7 would pollute that layer's
    query results.
\end{itemize}

The nine-layer decomposition is the smallest we found that keeps
each agent's write set disjoint across the most common concurrent
execution patterns.


%% ===================================================================
\section{The Ten Agents}\label{sec:agents}
%% ===================================================================

Table~\ref{tab:agents} summarises the ten agents. Each agent
inherits from a common \texttt{BaseAgent} class (181~lines) that
provides: Neo4j and Live-Memory client injection, a
\texttt{\_note()} dual-sink helper (Invariant~I1), an
\texttt{AgentExecution} recording wrapper (Invariant~I3), and
a typed \texttt{\_execute(**kwargs) -> dict} contract that each
agent overrides.

\begin{table*}[t]
\footnotesize
\centering
\caption{The ten agents: line counts, primary inputs, outputs, and
write layers. Line counts are for the agent module only, excluding
shared infrastructure.}
\label{tab:agents}
\begin{tabular}{@{}llrrll@{}}
\toprule
\textbf{Agent} & \textbf{Name} & \textbf{Lines} & \textbf{RFC} & \textbf{Input} & \textbf{Output}\\
\midrule
A1 & Ingestion              & 584 & \S5.1.1 & HldChangeset      & Intent, Device, VLAN\\
A2 & Intent$\rightarrow$Policy        & 221 & \S5.1.2 & Intent IDs        & Policy, FirewallRule\\
A3 & Policy$\rightarrow$Config        & 420 & \S5.1.2 & Neo4j context     & Configuration\\
A4 & Planning               & 243 & \S5.1.2 & Changeset+configs & MigrationPlan\\
A5 & Orchestration          & 350 & \S5.1.3 & Configuration     & DeploymentEvent\\
A6 & Monitoring             & 380 & \S5.2.1 & Device list       & Telemetry, OpState\\
A7 & Assessment             & 417 & \S5.2.2 & POR vs As-Built   & Assessment, Incident\\
A8 & Action                 & 405 & \S5.2.3 & Incident          & Remediation\\
A9 & Abstraction            & --- & \S5.2   & Assessments       & Aggregated scores\\
A10& Reporting              & --- & \S5.2   & Aggregations      & Reports\\
\bottomrule
\end{tabular}
\end{table*}

The total implementation of the eight active agents (A1--A8) is
3,020~lines of Python. Adding A9 and A10 (currently stubs) will
bring the total to approximately 3,500~lines.

\paragraph{Dispatch-table architecture.}
A recurring pattern across agents is the use of dispatch tables to
decouple vendor- or type-specific logic from the agent's control
flow. Three dispatch tables are central:

\begin{enumerate}
  \item \textbf{\texttt{\_KIND\_MAP}} (Provisioner): maps
    \texttt{(vendor, platform)} to containerlab kind, image, and
    type. Adding Nokia SR~Linux, FRR, and VyOS required one entry
    each---three lines total.
  \item \textbf{\texttt{\_TEMPLATE\_CHAINS}} (Agent~3): maps
    platform to an ordered list of Jinja2 template filenames. Adding
    FRR required one entry (three templates); adding SR~Linux
    required one entry (three templates).
  \item \textbf{\texttt{\_VENDOR\_EXECUTORS}} (Agent~5): maps
    platform to a callable that pushes configuration. Adding the
    gNMI transport for SR~Linux required one new callable and a
    two-line dispatch override keyed on
    \texttt{IBN\_SRLINUX\_TRANSPORT}.
\end{enumerate}

This architecture-as-data pattern means that adding a fourth vendor
(e.g.\ Arista cEOS) requires: one \texttt{\_KIND\_MAP} entry, one
\texttt{\_TEMPLATE\_CHAINS} entry (pointing to new Jinja2
templates), one executor function, and a Dockerfile for the container
image. No existing agent code changes.

The detailed operation of each agent is described in
Section~\ref{sec:cl:pipeline}. Here we note one additional design
decision per agent that is not covered by the pipeline description:

\begin{description}
  \item[A1:] Interactive dialogue support. When a population VLAN
    is ambiguous (e.g.\ multiple valid access-list templates apply),
    Agent~1 can invoke a \texttt{dialogue\_fn} callback with a
    prompt and a list of choices (\texttt{\_ACCESS\_TEMPLATES}: four
    templates labelled a/b/c/d). In batch mode the dialogue is
    auto-resolved to the first template; in interactive mode the
    operator selects.

  \item[A2:] Zone inference is rule-based, not LLM-based. The
    mapping from target string to source zone
    (\texttt{DMZ}~$\rightarrow$~\texttt{DMZ\_INFRA},
    \texttt{guest}~$\rightarrow$~\texttt{GUEST}, else
    \texttt{USER}) is three lines of code. This is deliberately
    minimal: the HLD's zone assignment table is authoritative, and
    the agent's role is to decompose, not to infer.

  \item[A3:] Content-hash-based change detection. Each rendered
    configuration is SHA-256 hashed (truncated to 16 hex chars).
    If the hash matches the previous configuration for the same
    device, the render is considered a no-op and no new
    \texttt{Configuration} node is written. This prevents redundant
    deployment events when the HLD is committed with cosmetic
    changes that do not affect device state.

  \item[A4:] Supersession logic. If a \texttt{PENDING} migration
    plan already exists when a new commit is processed, the old
    plan is marked \texttt{SUPERSEDED} rather than deleted. The
    superseded plan stays in the graph for audit, linked to the
    new plan via a \texttt{PRECEDED\_BY} edge.

  \item[A7:] Decorator-based drift detector registration. New
    drift detectors are added by writing a function decorated with
    \texttt{@\_drift\_detector("name")}; the decorator registers it
    in a global \texttt{\_DRIFT\_DETECTORS} dict. The assessment
    loop iterates over all registered detectors. Three are currently
    registered: \texttt{firewall\_rule}, \texttt{interface\_state},
    \texttt{vrrp\_role}.

  \item[A8:] Severity classification is intentionally decoupled
    from Agent~7's. Agent~7 classifies assessment severity (how bad
    is the drift?); Agent~8 classifies action severity (what should
    we do about it?). The two classifications use different inputs:
    A7 uses drift-event count and compliance percentage; A8 uses
    business-impact signals (service impact, redundancy status,
    device count).
\end{description}


%% ===================================================================
\section{Three-Tier Memory}\label{sec:memory}
%% ===================================================================

\subsection{Architecture Overview}

The three-tier memory architecture separates three qualitatively
different kinds of knowledge:

\begin{description}
  \item[Tier~1: Neo4j SSoT (\emph{what is}).]
    Structured, typed, queryable network state across nine layers.
    Accessed via Bolt (port 7687). Every agent reads and writes
    this tier on every cycle. Latency: $\sim$200\,ms per Cypher
    query.

  \item[Tier~2: Live-Memory (\emph{why it is so}).]
    Working memory for agent reasoning. Exposed via Model Context
    Protocol (MCP) Streamable HTTP on port 8002, backed by S3
    object storage. Each model state has a dedicated space (e.g.\
    \texttt{ibn-candidate-001}, \texttt{ibn-por-v1},
    \texttt{ibn-deploy-evt1}, \texttt{ibn-asbuilt-snap1}) plus
    two loop-scoped spaces (\texttt{ibn-loop-inner},
    \texttt{ibn-loop-outer}). Agents write structured notes via
    \texttt{live\_note()} calls that include: agent ID, timestamp,
    operation type, inputs summary, outputs summary, and
    free-text rationale. The notes accumulate in a \emph{bank}
    per space; the bank is periodically consolidated.

  \item[Tier~3: Graph-Memory (\emph{what was learned}).]
    Long-term knowledge graph with vector embeddings. Exposed via
    MCP on the same port (namespaced as \texttt{IBN\_LIFECYCLE\_*}
    labels in a separate Neo4j instance) plus Qdrant for BGE-M3
    1024-dimension embeddings. This tier stores entities and
    relations extracted from consolidated bank files using an
    LLM guided by the IBN Lifecycle Ontology. Agents query this
    tier via \texttt{question\_answer()} for historical
    precedent-based decision making (graph-guided RAG).
\end{description}


\subsection{The IBN Lifecycle Ontology}\label{sec:memory:ontology}

The Graph-Memory entity extraction is guided by a domain-specific
ontology (\texttt{ibn\_lifecycle\_ontology.yaml}, 472~lines) that
defines 14 entity classes across four families and 8 relationship
types.

\paragraph{Entity families.}

\begin{enumerate}
  \item \textbf{Decisions \& Rationale} (5 classes):
    \texttt{PolicyDecision} (firewall/zone decisions),
    \texttt{DesignChoice} (architectural decisions),
    \texttt{TradeOff} (competing approaches evaluated),
    \texttt{RejectedAlternative} (explicitly rejected options),
    \texttt{ApprovalRecord} (human sign-offs with operator identity
    and timestamp).

  \item \textbf{Incidents \& Remediation} (5 classes):
    \texttt{DriftEvent} (POR vs.\ As-Built mismatches),
    \texttt{RootCauseAnalysis} (causal chains for drift),
    \texttt{RemediationAction} (auto-fix or rollback records),
    \texttt{EscalationRecord} (human escalation events),
    \texttt{ComplianceTrend} (compliance trajectory across
    iterations: improving, stable, degrading).

  \item \textbf{Operations \& Deployment} (4 classes):
    \texttt{DeploymentNarrative} (config push records with
    per-device outcome),
    \texttt{ConfigurationRationale} (why a template rendered the
    way it did---e.g.\ ``VLAN~100 assigned to subinterface~100
    because the HLD population table maps Employee to VLAN~100''),
    \texttt{RollbackRecord} (rollback events with before/after
    config hashes),
    \texttt{VerificationResult} (post-push verification outcomes).

  \item \textbf{Knowledge \& Patterns} (4 classes):
    \texttt{RecurringPattern} (drift/latency/error correlations
    observed across multiple cycles),
    \texttt{LessonLearned} (operational guidance distilled from
    incident history),
    \texttt{CapacityInsight} (resource constraints---e.g.\ ``host
    memory gates provisioning beyond 12 nodes''),
    \texttt{CorrelationFinding} (metric correlations discovered
    by the aggregation pipeline).
\end{enumerate}

\paragraph{Relationship types.}
Eight typed relationships connect entities across families:

\begin{enumerate}
  \item \texttt{CAUSED\_BY}: drift/remediation to root cause.
  \item \texttt{RESOLVED\_BY}: incident to fix action.
  \item \texttt{PRECEDED}: temporal ordering of events.
  \item \texttt{JUSTIFIED\_BY}: decision to supporting rationale.
  \item \texttt{CONTRADICTS}: new finding vs.\ old lesson/pattern.
  \item \texttt{CONFIRMS}: event that validates a known pattern.
  \item \texttt{SIMILAR\_TO}: cross-iteration pattern matching
    (primary RAG retrieval link).
  \item \texttt{INFORMED\_BY}: decision informed by historical
    precedent (explicit query provenance).
\end{enumerate}

\paragraph{Extraction hints.}
The ontology includes natural-language extraction hints (lines
432--472) that guide the LLM's entity recognition. For example,
the \texttt{DriftEvent} class is triggered by phrases like
``drift detected'', ``non-compliant'', ``missing rule''; the
\texttt{LessonLearned} class by ``lesson:'', ``best practice'',
``should always''. These hints are injected into the LLM's
system prompt during \texttt{graph\_push()} to improve extraction
precision without fine-tuning.


\subsection{Consolidation Pipeline}\label{sec:memory:consolidation}

The \texttt{ConsolidationManager} (395~lines) orchestrates the
Tier~2$\rightarrow$Tier~3 data flow:

\begin{enumerate}
  \item \textbf{Trigger.} Consolidation fires on three events:
    (a)~a model-state transition
    (e.g.\ \texttt{CANDIDATE}$\rightarrow$\texttt{POR} after
    approval), (b)~the inner loop completing a full cycle, or
    (c)~the note count in a Live-Memory space exceeding a
    configurable threshold (default: 20 notes).

  \item \textbf{Consolidate.} The manager calls
    \texttt{bank\_consolidate(space\_id)} on the Live-Memory
    MCP server, which invokes an LLM to summarise the accumulated
    notes into a structured bank file. The LLM is configured via
    \texttt{LLMAAS\_API\_URL} and \texttt{LLMAAS\_MODEL}; the
    current testbed defaults to Claude Haiku~4.5 via Cloud-Temple's
    LLMaaS proxy, with Qwen3-235B as an alternative. The
    consolidation is space-specific: each model-state space has
    its own consolidation cadence and rules.

  \item \textbf{Push.} If the transition is \emph{push-eligible}
    (to \texttt{POR}, \texttt{DEPLOYED}, or \texttt{AS\_BUILT}),
    the manager reads all bank files via
    \texttt{bank\_read\_all(space\_id)} and pushes them to
    Graph-Memory via \texttt{graph\_push\_batch(memory,
    bank\_files, space\_id)}. The Graph-Memory server base64-encodes
    each bank file, invokes the LLM with the IBN Lifecycle Ontology
    as context, extracts typed entities and relations, embeds them
    with BGE-M3 (1024 dimensions), and stores them in Qdrant +
    the Graph-Memory Neo4j instance (namespaced under
    \texttt{IBN\_LIFECYCLE\_*} labels).

  \item \textbf{Query.} Any agent can query Tier~3 via
    \texttt{question\_answer(memory, question, max\_results=5)},
    which performs graph-guided RAG: the question is embedded,
    nearest-neighbour entities are retrieved from Qdrant, their
    graph neighbourhood is traversed via \texttt{SIMILAR\_TO} and
    \texttt{INFORMED\_BY} edges, and the combined context is
    returned to the caller.
\end{enumerate}

The consolidation pipeline is asynchronous and non-blocking: it
runs after the pipeline cycle completes and never delays the
fulfilment or assurance path. The manager also runs a background
sweep loop (default interval: 3,600\,s) that checks all spaces
for threshold-exceeding note counts, providing a safety net for
spaces that miss event-driven triggers.


\subsection{LLM Backend Configuration}\label{sec:memory:llm}

The LLM is used in exactly two places: Live-Memory consolidation
(bank summarisation) and Graph-Memory ingestion (entity/relation
extraction). Neither is in the request path of any pipeline stage.

The backend is configurable at deploy time via environment
variables:

\begin{lstlisting}[language=bash,numbers=none,caption={LLM configuration in \texttt{docker-compose.yml}.}]
LLMAAS_API_URL=http://host.docker.internal:8004/v1
LLMAAS_API_KEY=${LLMAAS_API_KEY}
LLMAAS_MODEL=${LLMAAS_MODEL:-claude-haiku-4-5-20251001}
LLMAAS_EMBEDDING_MODEL=all-MiniLM-L6-v2
LLMAAS_EMBEDDING_DIMENSIONS=1024
\end{lstlisting}

Swapping or disabling the LLM backend degrades long-term memory
quality (fewer entities extracted, coarser summaries) but does not
affect the pipeline's ability to parse, render, push, and verify.
The deterministic pipeline and the LLM-assisted memory are
independently testable and independently deployable.


%% ===================================================================
\section{Implementation}\label{sec:impl}
%% ===================================================================

\subsection{Testbed}\label{sec:impl:testbed}

The prototype runs on a single 12\,GB ARM64 virtual machine (Ubuntu
24.04, Linux 7.0.0, Python 3.14) with Docker and
Containerlab~\cite{containerlab}. The lab topology
(\texttt{clab-ibnlab-switches.yml}) defines 16 containers:

\begin{itemize}
  \item 3 Nokia SR~Linux (\texttt{ghcr.io/nokia/srlinux:latest},
    kind \texttt{nokia\_srlinux}, IXR-D2L profile): \texttt{acc-sw-01}
    (pioneer from Slice~2), \texttt{acc-01}, \texttt{agg-01}.
  \item 2 FRRouting (\texttt{ibn-frr:local}, custom Debian-slim
    image, 62\,MB, multi-arch): \texttt{edge-01}, \texttt{edge-02}.
  \item 4 VyOS firewalls (\texttt{ghcr.io/sysoleg/vyos-container:latest}):
    \texttt{usf-01}, \texttt{usf-02}, \texttt{dmzfw-01},
    \texttt{dmzfw-02}.
  \item 7 host stubs (lightweight Linux containers simulating end
    hosts in each population VLAN).
\end{itemize}

The management network is \texttt{192.168.100.0/24}. SR~Linux
containers expose gNMI on port 57400 and NETCONF on port 830 by
default. FRR containers run \texttt{zebra}, \texttt{bgpd},
\texttt{ospfd}, \texttt{ospf6d}, and \texttt{isisd} daemons with
\texttt{tini} as PID~1. The FRR image is built from a committed
Dockerfile (\texttt{infra/frr-image/Dockerfile}) to work around
the upstream \texttt{frrouting/frr} image being amd64-only.

Infrastructure services (Neo4j 5.x, Live-Memory MCP, Graph-Memory
MCP, Qdrant, Redis) run via \texttt{docker-compose.yml} on the
same host. Total memory usage with all 16 containers plus
infrastructure: $\sim$10\,GB, leaving $\sim$1.9\,GB headroom.


\subsection{Incremental Delivery: Six Slices}\label{sec:impl:slices}

The implementation was delivered in six incremental slices, each
building on the previous and each validated with unit tests and a
live end-to-end run. Table~\ref{tab:slices} summarises the slices.

\begin{table*}[t]
\footnotesize
\centering
\caption{Six implementation slices: scope, key deliverables, and
test counts.}
\label{tab:slices}
\begin{tabularx}{\textwidth}{@{}llXr@{}}
\toprule
\textbf{Slice} & \textbf{Name} & \textbf{Key deliverable} & \textbf{Tests}\\
\midrule
1 & HLD Pipeline &
  HLD parser, A1 ingestion, A2 decomposition, A3 rendering
  (VyOS only), A4 planning, A5 push via SSH, git hook trigger &
  52\\
2 & Multi-vendor &
  SR~Linux template chain (srl\_base, srl\_interfaces,
  srl\_vlans), docker exec sr\_cli executor,
  pioneer switch acc-sw-01 push &
  72\\
3 & Provisioning &
  Containerlab YAML mutation, provision() and
  retire(), \_KIND\_MAP dispatch, lifecycle events &
  86\\
4 & Approval gate &
  A4 severity classifier, MigrationPlan node, pending-plan
  Markdown, approve/reject CLI tool, resume\_from\_approval(),
  auto-approve for LOW/INFO &
  100\\
5 & FRR + batch &
  FRR as third vendor (Dockerfile, 3 templates, vtysh executor),
  provision\_batch() fix, 4 new devices provisioned,
  VM resize to 12\,GB &
  118\\
7a & gNMI transport &
  pygnmi executor for SR~Linux, CLI$\rightarrow$YANG
  translator (10 rules), IBN\_SRLINUX\_TRANSPORT env flag,
  live SetRequest + verification &
  140\\
\bottomrule
\end{tabularx}
\end{table*}

Each slice was committed with a plan document
(\texttt{PLAN\_Slice<N>.md}) before implementation and a fix-notes
document (\texttt{FIX\_NOTES\_Slice<N>.md}) after, capturing bugs
found, fixes applied, architectural decisions made during execution,
and known limitations. This convention was adopted after Slice~2 and
applied retroactively; it provides a readable post-mortem record
that survives outside the development conversation log.


\subsection{Failure Modes and Fixes}\label{sec:impl:failures}

Several non-trivial failures were encountered and resolved during
implementation:

\begin{description}
  \item[MCP hang on Python 3.14.]
    The \texttt{asyncio.run()} call hung in teardown after the MCP
    \texttt{streamablehttp\_client} context manager exited, due to
    an \texttt{anyio} TaskGroup orphan-task bug. After evaluating
    factory patterns, \texttt{asyncio.wait\_for}, and
    \texttt{ThreadPoolExecutor} (all of which also hung), the fix was
    a persistent background event loop in a daemon thread with
    \texttt{asyncio.run\_coroutine\_threadsafe()}. This pattern is
    now used by both \texttt{LiveMemoryClient} and
    \texttt{GraphMemoryClient}.

  \item[Batch provisioning timeout.]
    Per-device \texttt{clab deploy --reconfigure} re-processes the
    entire topology on each invocation. With four new devices, the
    third sequential deploy exceeded the 300\,s timeout. The fix
    (\texttt{provision\_batch()}) performs a single YAML write and
    a single deploy for the entire batch, completing in $\sim$10\,s.

  \item[FRR ARM64 mismatch.]
    The upstream \texttt{frrouting/frr:latest} image is amd64-only
    and crashes under QEMU on the ARM64 host. The fix was a custom
    \texttt{ibn-frr:local} image built from
    \texttt{debian:stable-slim} + \texttt{apt install frr}
    (62\,MB, native ARM64).

  \item[Neo4j OOM under memory pressure.]
    On the original 8\,GB host, Neo4j was OOM-killed (exit~137)
    when the fifth SR~Linux container came up alongside four VyOS
    firewalls. The host was resized to 12\,GB, which provides
    $\sim$4\,GB free after all containers and infrastructure.

  \item[Graph-Memory MCP schema mismatch.]
    The \texttt{GraphMemoryClient} initially sent field names that
    did not match the MCP server's expected schema (\texttt{name}
    instead of \texttt{memory\_id}, \texttt{content} instead of
    \texttt{content\_base64}). The \texttt{LLMAAS\_API\_URL} was
    also missing the \texttt{/v1} path prefix. Both were fixed by
    reading the server's OpenAPI spec and aligning field names.
\end{description}


\subsection{Measurements}\label{sec:impl:measurements}

\begin{table}[t]
\small
\centering
\caption{Pipeline timing for representative operations.}
\label{tab:timing}
\begin{tabular}{@{}lr@{}}
\toprule
\textbf{Operation} & \textbf{Wall-clock time}\\
\midrule
No-op HLD commit (full cycle) & 8--10\,s\\
VLAN addition (1 device affected) & 10--12\,s\\
Device addition (4 devices, batch) & 20--40\,s\\
gNMI SetRequest (7 updates) & $<$1\,s\\
sr\_cli push (typical) & $\sim$2\,s\\
vtysh push (typical) & $\sim$1.5\,s\\
SSH push to VyOS & $\sim$3\,s\\
Neo4j Cypher query (avg of 10 A3 queries) & $\sim$200\,ms\\
\bottomrule
\end{tabular}
\end{table}


%% ===================================================================
\section{Discussion and Related Work}\label{sec:discuss}
%% ===================================================================

\subsection{Comparison with Vendor IBN Suites}

Cisco DNA Center, Juniper Apstra, and Arista CloudVision all
implement forms of closed-loop intent, but they share three
properties our system avoids:

\begin{enumerate}
  \item \textbf{Vendor lock-in.} Each system manages only its own
    vendor's devices. Our dispatch-table architecture supports three
    vendors (VyOS, SR~Linux, FRR) with a fourth (Arista cEOS)
    requiring only a new dispatch entry and template chain.

  \item \textbf{Proprietary intent model.} Intent is authored in
    a vendor-specific GUI or API; it cannot be diffed, branched,
    or code-reviewed with standard tooling. Our intent is a
    Markdown file under \texttt{git}.

  \item \textbf{Opaque state.} The internal state model is not
    directly queryable by operators. Our Neo4j SSoT with its
    nine-layer ontology is fully queryable via Cypher, and every
    agent write is dual-sinked to Live-Memory for rationale
    capture.
\end{enumerate}

\subsection{Relationship to MALT}

Google's MALT~\cite{malt} introduced the concept of modelling
network topologies at multiple levels of abstraction, driven by
the observation that no single abstraction level serves all
consumers. Our nine-layer ontology is directly inspired by MALT
but differs in scope: MALT focuses on topology abstraction for
traffic engineering in a homogeneous datacentre; our ontology spans
the full IBN lifecycle (from business intent through deployment
to compliance) in a heterogeneous multi-vendor campus. The five
concurrent model states extend MALT's abstraction-level concept
with temporal semantics: the same intent exists simultaneously as
\texttt{CANDIDATE}, \texttt{POR}, \texttt{DEPLOYED}, and
\texttt{AS\_BUILT}, each reflecting a different point in the
fulfilment lifecycle.

\subsection{GitOps for Networks}

The GitOps pattern---declarative desired state in a Git repository,
reconciled by an automated operator---is well-established for
Kubernetes~(Flux, ArgoCD) but less explored for network
infrastructure. Network-to-Code's Nautobot and Netbox provide
source-of-truth databases but do not close the loop: they record
intent but do not render, push, or verify. Our system closes the
loop end-to-end: the HLD commit triggers rendering, the approval
gate provides human oversight, and the assurance loop detects
drift and auto-remediates.

\subsection{YANG and NETCONF/gNMI}

The move from CLI-based config push to schema-validated gNMI
(Slice~7a) aligns with the broader industry trend toward
model-driven management~\cite{gnmi}. Our experience confirms
that gNMI's structured \texttt{SetRequest} semantics (atomic
commit, typed paths, JSON-IETF encoding) are a better fit for
programmatic push than piping CLI commands through
\texttt{docker exec}. The narrow CLI$\rightarrow$YANG translator
(\texttt{\_srlinux\_yang.py}) is a pragmatic compromise: it keeps
the A3 template chain transport-agnostic (the same \texttt{set /}
commands feed both the CLI and gNMI paths) while the translator
absorbs the encoding difference. A future slice (7c) will render
OpenConfig JSON directly, eliminating the translation step.

\subsection{Lessons Learned}

\begin{enumerate}
  \item \textbf{The HLD-as-intent commitment is viable but
    demanding.} The HLD must be structured enough for a
    deterministic parser to extract typed records. In practice
    this means the Device Inventory Table and VLAN assignment
    tables must follow a fixed Markdown table format. Prose
    sections (design rationale, capacity planning) are not parsed
    but are preserved in Live-Memory notes for future LLM-driven
    extraction.

  \item \textbf{Batch provisioning matters.} The per-device
    \texttt{clab deploy --reconfigure} path scales quadratically
    with topology size because each invocation re-processes the
    entire YAML. The batch path (one write, one deploy) reduces
    this to constant time for the deploy step.

  \item \textbf{Three tiers are the right number.} Early designs
    considered a two-tier model (Neo4j + LLM-backed knowledge).
    The working-memory tier (Live-Memory) turned out to be
    essential: agents need a fast, low-latency scratchpad for
    intra-cycle communication that is richer than Neo4j properties
    but cheaper than a full knowledge-graph ingestion.

  \item \textbf{LLMs belong in the memory tier, not the pipeline.}
    Every attempt to introduce LLM inference into the pipeline
    (e.g.\ LLM-based policy decomposition in A2, LLM-based
    template selection in A3) degraded cycle time by 10--30$\times$
    and introduced non-determinism that broke idempotency
    invariants. The current separation---deterministic pipeline,
    LLM-assisted memory---is a deliberate architectural choice,
    not a limitation.

  \item \textbf{Dispatch tables scale better than inheritance.}
    The vendor-extensibility pattern (one table entry per vendor)
    has proven more maintainable than a class-hierarchy approach.
    Adding FRR as a third vendor touched exactly one line per
    dispatch table plus a Dockerfile and three Jinja2 templates;
    no existing code was modified.
\end{enumerate}


%% ===================================================================
\section{Conclusion}\label{sec:conclusion}
%% ===================================================================

We have presented an Intent-Based Networking closed-loop control
system that treats a conventional Markdown High-Level Design
document as the sole source of intent, triggered by \texttt{git
commit} and reconciled to live network state through a pipeline of
ten specialised agents grounded in RFC~9315. The system is backed
by a Neo4j knowledge graph with a nine-layer ontology and five
concurrent model states, a three-tier memory architecture that
cleanly separates structured state, working memory, and long-term
knowledge, and a dispatch-table-driven multi-vendor rendering and
push pipeline supporting VyOS, Nokia SR~Linux (via both CLI and
gNMI), and FRRouting.

The implementation was delivered incrementally across six slices on
a live Containerlab testbed with 16 nodes, validated by 140 unit
tests with zero cross-slice regressions, and operates with a
full-cycle time of 8--10\,seconds. The core pipeline is entirely
deterministic and LLM-free; large language models are used
exclusively for memory consolidation and knowledge extraction,
preserving the system's correctness properties while enabling
rich operational knowledge capture.

The key architectural insight is that the gap between an HLD and a
running network is not a translation problem that requires a
proprietary controller---it is a reconciliation problem that can be
decomposed into narrow, testable agents operating over a shared
graph, triggered by the same version-control tooling that
developers already use. The HLD is not a specification of intent
that must be re-encoded into a machine format; it \emph{is} the
intent, and the closed loop's job is to make reality match it.

\paragraph{Future work.}
Immediate next steps include: NETCONF-over-SSH transport for
SR~Linux (Slice~7b), OpenConfig shared templates across vendors
(Slice~7c), gNMI streaming telemetry for sub-second inner-loop
monitoring (Slice~7d), and full implementation of Agents~9 and~10
for abstraction and human-facing reporting. Longer-term, we plan
to evaluate predictive drift detection using Graph-Memory's
historical knowledge base and to extend the testbed to include
Arista cEOS and Juniper cRPD.


%% --------- bibliography ---------
\bibliographystyle{ACM-Reference-Format}
\bibliography{refs}

\end{document}
